\section{Related Work}
\label{sec:related_work}
% \todo{Section dedicated to the related work, i.e., similar / relevant studies, and how your study fits in the overall research picture. Usually, this section includes secondary studies which address the same / similar topic to the one addressed in your study.}

This systematic literature review builds upon a growing body of research exploring the convergence of HPC and cloud computing. Several existing studies have investigated various aspects of this transition, providing valuable insights into the challenges, opportunities, and future directions of this field.

\noindent
\textbf{Early Research on HPC in the Cloud}
Early research focused on identifying the challenges of migrating HPC applications to cloud environments. Studies like \cite{wang2010cloud, gupta2014challenges, marathe2013challenges} highlighted the limitations of cloud networks, virtualization overhead, and the unpredictable nature of shared cloud resources. These studies emphasized the need for specialized tools and infrastructure to support HPC workloads in the cloud.

\noindent
\textbf{Recent Research on HPC-Cloud Convergence}
More recent research has explored the development of solutions and technologies to address the challenges identified in earlier studies. For example, \cite{netto_hpc_2019} provides a comprehensive overview of the state of the art in HPC-Cloud convergence, highlighting the importance of hybrid cloud models, containerization technologies, and resource disaggregation.
Several studies have focused on specific aspects of HPC-Cloud convergence, such as: \\
\newline
\noindent
\textit{Network Performance}: \cite{alves_interference-aware_2020, giannoutakis_evaluation_2021} investigate the performance of cloud networks for HPC workloads, exploring the use of high-speed interconnects and optimized communication protocols.\\

\noindent
\textit{Virtualization Technologies}: \cite{marques_harnessing_2021} examines the use of containerization and operating-system-level virtualization to reduce overhead and improve performance in cloud-based HPC environments.\\

\noindent
\textit{Resource Disaggregation}: \cite{michelogiannakis_case_2022, copik_software_2024} explore the potential of resource disaggregation to optimize hardware utilization and enhance flexibility in HPC-Cloud systems.\\

\noindent
\textit{Cloud Storage for HPC}: \cite{gadban_analyzing_2022, guidi_10_2021} investigate the use of cloud storage services like Amazon S3 for HPC workloads, highlighting the benefits and challenges of this approach.\\

\noindent
\textit{AI and Machine Learning}: \cite{chahal_isesa_2022, michalewicz_attempt_2023, silva_review_2024} explore the use of AI and ML for optimizing HPC workloads, managing cloud resources, and enabling new scientific applications.\\

This review builds upon the existing literature by providing a comprehensive overview of the current state of the art in HPC-Cloud convergence, focusing on recent research (2021-2024) and highlighting emerging trends. The review also emphasizes the importance of a hybrid approach, combining on-premise HPC with cloud resources, as a practical and beneficial path for achieving convergence. Furthermore, the review explores the potential of AI and ML to address key challenges and drive innovation in the field.
By synthesizing and analyzing the findings of previous studies, this review aims to provide a valuable resource for researchers, practitioners, and policymakers interested in the future of HPC-Cloud convergence.
